% =====================================================================
\section{Introduction}
\label{sec:intro}
% =====================================================================

This document specifies the requirements for CISC-8: an 8-bit CISC processor
with a Harvard architecture, 64 general-purpose registers accessed through a
sliding window of 8 registers, and a hardware \IIC{} controller. It is the
baseline for the RTL design.

\subsection{Conventions}

Every requirement has a unique identifier of the form \texttt{PREFIX-NNN} and
is verifiable. The following keywords define the weight of a requirement:

\begin{center}
\small
\begin{tabularx}{\linewidth}{@{}lX@{}}
  \toprule
  \textbf{Keyword} & \textbf{Meaning} \\
  \midrule
  \MUST{} / \MUSTNOT & Mandatory; failure to comply is a non-conformance. \\
  \SHOULD & Recommended; any deviation must be justified in the design documentation. \\
  \MAY & Optional; left to the implementer. \\
  \bottomrule
\end{tabularx}
\end{center}

Prefixes: \texttt{ARCH} (general architecture), \texttt{MEM} (memory),
\texttt{REG} (registers), \texttt{SFR} (special function registers),
\texttt{ISA} (instructions), \texttt{I2C} (\IIC{} controller),
\texttt{INT} (interrupts and reset), \texttt{IO} (interfaces),
\texttt{NF} (non-functional).

Hexadecimal numbers are written with the prefix \texttt{0x}, binary numbers
with \texttt{0b}. Bit 0 is the least significant bit (LSB).

